\scanpage{153}
\mspage{724}

\unsure{ment}, \deleted{\illegible{1}} or to make explicit in terms of \unsure{standard} ``elliptic functions''.

Without dwelling on such an explicitation, and returning \inserted{rather} to the corresponding embedding
\diagram{%
(\unsure{318})\qquad
\begin{tikzcd}[ampersand replacement=\&,column sep=large]
U \arrow[r,hook] \& \mathcal J^0_{U/S}
\end{tikzcd}}
and to the \unsure{marked} section
\[
(\unsure{319})\qquad e\in U(S)=H^0(S,\mathcal J^0_{U/S}),
\]
\inserted{we can use the tower of isogenies}
\deleted{we can use the canonical homomorphism}
\diagram{%
\deleted{(\unsure{314})}\qquad
\begin{tikzcd}[ampersand replacement=\&,column sep=large]
\deleted{H^0(S,\mathcal J^0_{U/S})} \arrow[r,phantom,"\deleted{\longrightarrow}"{description}] \& \deleted{H}
\end{tikzcd}}
\diagram{%
(\unsure{321})\qquad
\begin{tikzcd}[ampersand replacement=\&,column sep=large]
1 \arrow[r,dash] \& {}_n\mathcal J^0_{U/S} \arrow[r,hook] \& \mathcal J^0_{U/S} \arrow[r,"n.\mathrm{id}_{\mathcal J^0_{U/S}}"] \& \mathcal J^0_{U/S} \arrow[r,dash] \& 1
\end{tikzcd}}
(for varying $n$), i.e. the universal isogeny deduced by passage to $\varprojlim$:
\diagram{%
(\unsure{322})\qquad
\begin{tikzcd}[ampersand replacement=\&,column sep=large,row sep=large]
1 \arrow[r,dash] \& T_\infty(\mathcal J^0_{U/S}) \arrow[r,dash] \arrow[d,phantom,"\rotatebox{90}{$\simeq$}"{description}] \& \widetilde{\mathcal J}^0_{U/S} \arrow[r,"p"] \arrow[d,phantom,"\rotatebox{90}{$\underset{\mathrm{def}}{=}$}"{description}] \& \mathcal J^0_{U/S} \arrow[r] \& 1 \\
\& T_\infty(\mathbb G_m)_S\otimes_{\mathbb Z}\unsure{M_0}(S'/S) \& \varprojlim_n\unsure{{}_n}(\mathcal J^0_{U/S}) \& \&
\end{tikzcd}}
to deduce, by inverse image of $e$, a torsor under $T_\infty(\mathcal J^0_{U/S})\simeq T_\infty(\mathbb G_m)_S\otimes_{\mathbb Z}\unsure{M_0}(S'/S)$.

This being so, we will have (cf. (301) page 719')
\[
(323)\qquad F_\infty(\widetilde X/S)\simeq p^{-1}(e),\quad\text{the }T_\infty(\mathcal J^0_{U/S})\text{-torsor}
\]
\inserted{deduced from the section $e$ of $\mathcal J^0_{U/S}$ over $S$.}
\margin{152}

\scanpage{154}

\uline{Remark} The isomorphism class of the $T_\infty(\mathcal J^0_{U/S})$-torsor obtained is an element of $H^1(S,T_\infty(\mathcal J^0_{U/S}))$, image of the element $e$ of $H^0(S,\mathcal J^0_{U/S})$, under the canonical morphism
\diagram{%
(324)\qquad
\begin{tikzcd}[ampersand replacement=\&,column sep=large]
H^0(S,\mathcal J^0) \arrow[r] \& H^1(S,T_\infty(\mathcal J^0))
\end{tikzcd},}
whose kernel consists of the subgroup of $H^0(S,\mathcal J^{\unsure{0}})$ of ``infinitely divisible'' elements. In the case where $S$ is \inserted{\unsure{reduced and is}} of finite type over \deleted{\unsure{$\mathbb Q$}} $\mathbb Z$ -- for example \inserted{reduced} of finite type over $\mathbb Q$ -- we know (and this is \deleted{\unsure{a theorem}} ``elementary'') that this kernel is reduced to zero. This therefore means that the section $e$ \inserted{of $\unsure{U}$ over $S$} (i.e. the ``elliptic parameter'' under consideration) is then entirely determined, when one knows
\margin{153}

\scanpage{155}
\mspage{725}

$1^{\circ})$ the \deleted{\illegible{1}} associated étale triple covering $S'$ of $S$, corresponding to the points of order $2$ of $\widetilde X/S$ -- whence $\mathcal J^0\unsure{=}\mathcal J^0_{U/S}\simeq\mathcal J^0_0{}_{\deleted{\unsure{S}}}(\widetilde X,S')$ -- and

$2^{\circ})$ The $T_\infty(\mathcal J^{\unsure{0}}_{\deleted{\unsure{U/S}}})$-torsor $\unsure{\mathbb F_\infty|F_\infty}(\widetilde X/S)$;
\margin{154}

\scanpage{156}

$15^{\circ})$ Interpretation of $\Gamma_{\mathbb Q}\longrightarrow N_\sigma^{\unsure{\natural|\sharp}}$.

In number $11^{\circ}$ we introduced the canonical homomorphism
\diagram{%
\begin{tikzcd}[ampersand replacement=\&,column sep=huge,row sep=large]
\mathcal M\overset{\mathrm{def}}{=}\mathcal N_{1,1}^{\text{ét}} \arrow[r] \arrow[d,phantom,"\cup"{description}] \& \mathcal M_{\rho,\sigma}^{\sim} \arrow[d,phantom,"\cup"{description}] \\
\mathcal M^! \arrow[r] \& \mathcal M_{\rho,\sigma}^{\sim\unsure{\natural|\sharp}}
\end{tikzcd}\qquad\text{cf. (217) page \unsure{695})}}
which we have ``decomposed'' into \deleted{more than} two homomorphisms (in addition to a trivial homomorphism
\diagram{%
\begin{tikzcd}[ampersand replacement=\&,column sep=large,row sep=large]
\mathcal M \arrow[r] \arrow[dr] \& N_\rho^{\sim}\simeq\widehat{\mathbb Z}^{\times} \\
\& \Gamma_{\mathbb Q}^{\text{ét}} \arrow[u,"\chi"']
\end{tikzcd}}
) into \uline{two} homomorphisms
\diagram{%
(\unsure{325})\qquad
\begin{tikzcd}[ampersand replacement=\&,column sep=huge,row sep=large]
\mathcal M \arrow[r] \arrow[d,phantom,"\cup"{description}] \& G_{\rho,\sigma}\simeq\widehat{\mathbb Z}^{*}.\mathfrak S^{+\wedge} \arrow[d,phantom,"\unsure{\cup}"{description}] \\
\mathcal M^! \arrow[r] \& G_{\rho,\sigma}^{\unsure{\natural|\sharp}}\simeq\widehat{\mathbb Z}^{*}.M^\wedge
\end{tikzcd}\qquad\text{(cf. (218))}}
and
\diagram{%
(\unsure{326})\qquad
\begin{tikzcd}[ampersand replacement=\&,column sep=large]
\Gamma_{\mathbb Q}^{\text{ét}} \arrow[r] \& N_\sigma^{\unsure{\natural|\sharp}\deleted{\unsure{\mathrm{ab}|\sim}}}\simeq\widehat{\mathbb Z}^{*}.\widehat{\mathbb Z}
\end{tikzcd}\qquad\text{(cf. (219)).}}

The first homomorphism (\unsure{322}) was characterized in \inserted{\unsure{arithmetico-}}\deleted{\unsure{algebraic}}geometric terms (as regards its restriction to $\mathcal N_{1,1}\subset\mathcal N_{1,1}^{\text{ét}}$ itself, and reducing modulo $\mu_G$ to find a homomorphism
\diagram{%
(327)\qquad\deleted{\unsure{$\mathcal M_{0,3}$}}\qquad
\begin{tikzcd}[ampersand replacement=\&,column sep=large]
\mathcal M_{0,3} \arrow[r,dash] \& G_{\rho,\sigma}^{\sim}\overset{\mathrm{def}}{\simeq}G_{\rho,\sigma}/\mu_G
\end{tikzcd},}
in section $13^{\circ})$, via the tower of
\margin{155}

\scanpage{157}
\mspage{\unsure{725|726}}

Fermat curves. From the present point of view, where we are trying to make explicit the algebraic geometry of \inserted{\unsure{essential\textsuperscript{ly}}} finite type over $\mathbb Q$ in terms of \unsure{automorphisms} of profinite groups, we can also say, conversely, that we have managed to ``describe'' the tower of Fermat curves, via the homomorphism (327) constructed here ``by brute force''.

It would now remain to find an arithmetico-geometric explicitation of the homomorphism \deleted{\unsure{(322)}} from $\Gamma_{\mathbb Q}$ \inserted{$\subset\Gamma_{\mathbb Q}^{\text{ét}}$} into $\widehat{\mathbb Z}^{*}.\widehat{\mathbb Z}$, induced by (\unsure{326}).
But \deleted{\unsure{this the homomorphism}} using the embedding
(\diagram{%
\begin{tikzcd}[ampersand replacement=\&,column sep=large]
N_\sigma^{\unsure{\natural|\sharp}\deleted{\unsure{\mathrm{ab}|\sim}}} \arrow[r,hook] \& \unsure{G}_{\rho,\sigma}^{\unsure{\natural|\sharp}\deleted{\unsure{\mathrm{ab}|\sim}}}\simeq\widehat{\mathbb Z}^{*}.\uline{M}
\end{tikzcd},}
the homomorphism (\unsure{313|323}) appears as the composite
\diagram{%
\begin{tikzcd}[ampersand replacement=\&,column sep=large]
\Gamma_{\mathbb Q}^{\text{ét}} \arrow[r,hook,"\unsure{\kappa}_{1/2}"'] \& \mathcal M^! \arrow[r,"(325)","\text{i.e. \textquotedblleft Fermat\textquotedblright}"'] \& G_{\rho,\sigma}^{\unsure{\natural|\sharp}\deleted{\unsure{\mathrm{ab}|\sim}}}
\end{tikzcd},}
more precisely we have a commutative diagram:
\margin{156}

\scanpage{158}
\[
\tag{328}
\diagram{\begin{tikzcd}[ampersand replacement=\&,column sep=large,row sep=large]
[\Gamma_{\mathbb Q}\hookrightarrow]\Gamma_{\mathbb Q}^{\text{ét}}
 \arrow[r] \arrow[d,hook,"\kappa_{1/2}"'] \&
 N_{\sigma}^{\unsure{\natural|\sharp}}\simeq\widehat{\mathbb Z}^{*}.\widehat{\mathbb Z}
 \arrow[d,hook]\\
{}[\underbrace{\mathcal M_{0,3}^{!}}_{\Pi_1(\mathcal U_{0,3\mathbb Q})}\hookrightarrow]
 \mathcal M^{!\,\unsure{\infty|\text{ét}}}\simeq\unsure{N|\mathcal N}_{\unsure{\rho\beta|\rho3}}^{!\,\text{ét}}
 \arrow[r,"{\substack{\text{``Fermat''}\\(325)}}"'] \&
 G_{\rho,\sigma}^{\unsure{\natural|\sharp}}\simeq\widehat{\mathbb Z}^{*}.\uline{M}^{\wedge}.
\end{tikzcd}}
\]

This therefore means that we find the homomorphism induced by (326)
\[
\tag{329}
\diagram{\begin{tikzcd}[ampersand replacement=\&]
\Gamma_{\mathbb Q}\arrow[r]\&N_{\sigma}^{\unsure{\natural|\sharp}}=\widehat{\mathbb Z}^{*}.\widehat{\mathbb Z}
\end{tikzcd}}
\]
by looking at the image \inserted{in $N_{\sigma}^{\unsure{\natural|\sharp}}\subset\widehat{\mathbb Z}^{*}.M^{\wedge}$} of the homomorphism
\[
\tag{330}
\diagram{\begin{tikzcd}[ampersand replacement=\&]
\Gamma_{\mathbb Q}\arrow[r]\&\widehat{\mathbb Z}^{*}.\uline{M}^{\wedge}\subset\operatorname{Aff}(M^{\wedge}),
\end{tikzcd}}
\]
\deleted{in $N_{\sigma}^{\unsure{\natural|\sharp}}\subset\widehat{\mathbb Z}^{*}.\uline M$,}
deduced from the tower of fermat curves by looking at the fiber of the said tower at the point $1/2$, which gives a torsor under
$T_{\infty}(\mathbb G_m)\otimes_{\widehat{\mathbb Z}}\widehat M$,
being made explicit Galois-theoretically precisely by the homomorphism (326).
Let us note that the \deleted{\unsure{fiber}}\inserted{point $1/2$}
\deleted{$U_{\infty}(1/2)$ is \unsure{canonically isomorphic},}
\inserted{of $\mathcal U_{0,3\mathbb Q}$ corresponds, under the isomorphism $\psi$ (317)}
of $\mathcal U_{0,3\mathbb Q}$ with $U_{1,\mathbb Q}$, to the point $(-\tfrac12,-\tfrac12,\unsure{1})$
\margin{157}

\scanpage{159}
\mspage{727}
of \deleted{\unsure{$U$}}
\[
\tag{331}
\inserted{s_{1/2}=(-\tfrac12,-\tfrac12,1)\in}
\deleted{\unsure{U_{1/2}}}
\left\{\begin{gathered}
U_{1\mathbb Q}\subset\mathbb G_{m\mathbb Q}^{3}/\mathbb G_{m\mathbb Q}\subset\mathbb P_{\mathbb Q}^{2}\\
\Vert\\
\{(x_0,x_1,x_{\infty})\in\Gamma(\mathbb G_{m\mathbb Q}^{3}/\mathbb G_{m\mathbb Q})\mid x_0+x_1+x_{\infty}=0\}
\end{gathered}\right.
\]
The \inserted{\unsure{geometric|geometric}} fiber of $U_n$ at the said point is thus identified with the systems of solutions of the equations below:
\[
\tag{332}
\begin{aligned}
\deleted{\unsure{\bar s|s}}\quad U_n(s_{1/2})
&\simeq\{(x_0,x_1,x_{\infty})\in\Gamma\mathbb G_m^3\deleted{/\mathbb G_m}\mid x_0^n=-\tfrac12,\ x_1^n=-\tfrac12,\ x_{\infty}^{n}=\unsure{1|-1}\}\\[-2pt]
&\hspace{5em}\deleted{\unsure{\mathrm{mod}}\ \mathbb G_m}\\
&\simeq\{(\unsure{x_0},x_1,x_{\infty})\in\Gamma(\mathbb G_m^3)\mid x_0=-\tfrac12,\ x_1^n=-\tfrac12,\ x_{\infty}^n=\unsure{1|-1}\}\\
&\simeq\{(x_0,x_{\unsure{1|\infty}})\in\Gamma\mathbb G_m^2\mid x_0^n=\unsure{-\tfrac12},\ x_{\unsure{1|\infty}}^n=-\tfrac12\}\\
&\simeq\deleted{\unsure{\mu_n\times}}T_n(-\tfrac12)\times T_n(\unsure{\tfrac12|-\tfrac12})
\end{aligned}
\]
where we have put, for $\forall\ \unsure{\zeta|\xi}\in\mathbb G_m(S)$ ($S$ any scheme, \inserted{here $S=\operatorname{Spec}\mathbb Q$})
\[
\tag{333}
\begin{gathered}
T_n(\unsure{\zeta|\xi})=
\deleted{\unsure{\text{image of}}\ \unsure{\zeta|\xi}\in H^0(S,\mathbb G_m)\ \text{under}}\\
\inserted{\mu_n\text{-torsor of }\deleted{\unsure{\text{of}}}n^{\text{th}}\text{ roots of }\unsure{\zeta|\xi}},
\end{gathered}
\]

Passing to the limit we find an isomorphism
\[
\tag{334}
U_{\infty}(s_{1/2})\simeq T_{\infty}(\unsure{-\tfrac12})\times T_{\infty}(-\tfrac12),
\]
where we put, as is only right
\[
\tag{335}
T_{\infty}(\unsure{\zeta|\xi})=\varprojlim_n T_n(\unsure{\zeta|\xi})
\qquad(\text{a }T_{\infty}(\mathbb G_m)\text{-torsor}).
\]
\margin{158}

\scanpage{160}
Thus, the homomorphism (329) should be interpretable in terms of the Galois homomorphism
\[
\tag{336}
\diagram{\begin{tikzcd}[ampersand replacement=\&]
\Gamma_{\mathbb Q}\arrow[r]\&\operatorname{Gal}\bigl(\widetilde{\mathbb Q}[\sqrt[\infty]{-\tfrac12}]/\mathbb Q\bigr)
\end{tikzcd}}
\]
where $\widetilde{\mathbb Q}$ denotes the universal cyclotomic extension of $\mathbb Q$
\[
\tag{336}
\widetilde{\mathbb Q}\simeq\mathbb Q[\sqrt[\infty]{1}].
\]

To appreciate the ``size'' of the image of $\Gamma_{\mathbb Q}$ in $\widehat{\mathbb Z}^{*}.\widehat{\mathbb Z}\simeq N_{\sigma}^{\unsure{\natural|\sharp}}$, it is appropriate to make a digression of a general nature, on the Galois groups \inserted{over $\mathbb Q$} of extensions of the form
\[
\widetilde{\mathbb Q}[\sqrt[\infty]{\unsure{\zeta_1|\xi_1}},\ldots,\sqrt[\infty]{\unsure{\zeta_N|\xi_N}}].
\]

For this, over a base $S/\mathbb Q$, \unsure{we have for simplicity} (we \unsure{will be|restrict ourselves} here \unsure{reduced|restricted} \unsure{to the} case $S=\operatorname{Spec}\mathbb Q$) let us give ourselves an abelian group $A$ (\inserted{in} the case that interests us here we will take $A=\mathbb Z$), and
\margin{159}

\scanpage{161}
\mspage{728}
a homomorphism
\[
\tag{337}
\diagram{\begin{tikzcd}[ampersand replacement=\&]
A\arrow[r,"\unsure{\zeta|\xi}"]\&\mathbb G_m(S)=H^0(S,\mathbb G_{mS})
\end{tikzcd}}
\]
(in the present case, it is given by $1\mapsto-1/2$),
or what amounts to the same thing, a homomorphism
\[
\tag{338\,\deleted{\unsure{bis}}}
\diagram{\begin{tikzcd}[ampersand replacement=\&]
A_S\arrow[r,"\unsure{\zeta|\xi}"]\&\mathbb G_{mS}.
\end{tikzcd}}
\]

Let us consider the tower of Kummer isogenies
\[
\tag{33\deleted{\unsure{8}}\inserted{\unsure{9}}}
\diagram{\begin{tikzcd}[ampersand replacement=\&]
1\arrow[r,dash]\&\mu_n\arrow[r,dash]\&\mathbb G_m\arrow[r,dash,"\lambda\mapsto\lambda^n"]\&\mathbb G_m\arrow[r,dash]\&1
\end{tikzcd}}
\]
whence, by inverse image via (33\unsure{8}), a tower of extensions of étale group schemes over $S$
\[
\tag{339}
\diagram{\begin{tikzcd}[ampersand replacement=\&]
1\arrow[r]\&\mu_n\arrow[r,dash]\&E_n(\unsure{\zeta|\xi})\arrow[r,dash]\&A_S\arrow[r]\&1
\end{tikzcd}}
\]
and by \unsure{passage|passage} to the $\varprojlim$
\[
\tag{340}
\diagram{\begin{tikzcd}[ampersand replacement=\&]
1\arrow[r,dash]\&T_{\infty}(\mathbb G_{mS})\arrow[r,dash]\&E_{\infty}(\unsure{\zeta|\xi})\arrow[r,dash]\&A_S\arrow[r]\&1.
\end{tikzcd}}
\]

Let us suppose $S$ is \inserted{$0$-}connected, and take a \deleted{\unsure{fixed}} geometric point \inserted{$\bar s$} of $S$, we will therefore have a canonical homomorphism
\[
\tag{341}
\diagram{\begin{tikzcd}[ampersand replacement=\&]
\Pi_1(S,\bar s)\arrow[r]\&\operatorname{Aut}_{\substack{\mathrm{gr}\\\text{profinite}}}\bigl(E_{\infty}(\unsure{\zeta|\xi})_{\bar s}\bigr),
\end{tikzcd}}
\]
\margin{160}